%% kapa: exam=1 round=30 year=2019 period=2 form=A subject=appraisal_law subject_ko=감정평가관계법규 seq=6 number=6
%% source: Q-Net 감정평가사 기출문제 (cst00309) artlSeq=5205205 file=extracted/1008/5205205/2019년 제30회 감정평가사 1차시험 문제지/2019년 제30회 감정평가사 1차시험 문제지/2019년도 제30회 감정평가사 1차 2교시 A형.hwp
%% conversion: kordoc 4.24.51 -> tools/gen1.py (auto); review=proofread-sonnet; flags=none
\begin{question}{6}
국토의 계획 및 이용에 관한 법령상 지구단위계획에 관한 설명으로 옳지 않은 것은?
\begin{choices}
  \item 지구단위계획구역의 지정에 관한 고시일부터 5년 이내에 지구단위계획이 결정ㆍ고시되지 아니하면 그 5년이 되는 날에 지구단위계획구역의 지정에 관한 도시ㆍ군관리계획결정은 효력을 잃는다.
  \item 지구단위계획에는 건축물의 건축선에 관한 계획이 포함될 수 있다.
  \item 지구단위계획구역 및 지구단위계획은 도시ㆍ군관리계획으로 결정한다.
  \item 국토교통부장관, 시ㆍ도지사, 시장 또는 군수는 지구단위계획구역 지정이 효력을 잃으면 지체 없이 그 사실을 고시하여야 한다.
  \item 국토교통부장관은 용도지구의 전부 또는 일부에 대하여 지구단위계획구역을 지정할 수 있다.
\end{choices}
\end{question}
\answer{1}
